\section{Introduction}
\label{sec:introduction}

Large language models (LLMs) are increasingly deployed as autonomous
agents capable of interacting with external tools and
APIs~\citep{schick2024toolformer, yao2023react}. Rather than generating
text in isolation, these agents retrieve documents, query databases,
send messages, and trigger downstream workflows---actions with
real-world side effects. A growing body of work surveys this
paradigm~\citep{wang2024agent_survey}, and several orchestration
frameworks have emerged to structure agent behavior as explicit workflow
graphs: LangGraph~\citep{langgraph_docs},
CrewAI~\citep{crewai_docs}, AutoGen~\citep{wu2023autogen}, and Google
ADK~\citep{google_adk_docs} each represent computation steps as nodes
and allowable transitions as edges.

Despite this rich structural representation, safety enforcement in
current agent systems remains predominantly a \emph{runtime} concern.
Tools such as NeMo Guardrails~\citep{nemo_guardrails} and
LlamaGuard~\citep{llamaguard} intercept individual LLM calls or tool
invocations and apply content-level or policy-level filters at execution
time. While valuable for catching toxic outputs and prompt injection
attempts, these runtime approaches have three fundamental limitations:
(i)~they impose per-call latency overhead, (ii)~they detect violations
only when the offending execution path is actually exercised, and
(iii)~they provide no coverage guarantees over paths that remain
untriggered during testing.

\paragraph{Motivating example.}
Consider an email triage agent with the following workflow: incoming
emails are classified by intent, routed to either an urgent or normal
handler, drafted, and sent. During iterative development, a developer
adds a \texttt{draft\_response} node on the normal-priority path but
forgets to connect it to the \texttt{send} node. The result is a
\emph{dead end}: normal-priority emails enter the draft stage and
silently halt. A runtime guardrail would not catch this defect unless
the normal-priority path is exercised with an email that triggers it.
A static analysis of the workflow graph, however, immediately flags the
dead-end node and produces a witness trace:
\texttt{\_\_start\_\_ $\to$ classify $\to$ router $\to$
normal\_handler $\to$ draft\_response} (stuck).

\paragraph{Research gap.}
This class of defect---a topological error in the workflow
graph---cannot be addressed by content-level runtime guardrails. In
principle, classical model checking techniques could detect such
errors: reachability analysis, deadlock detection, and temporal logic
verification are well-established in the formal methods
literature~\citep{clarke1999modelchecking, baier2008modelchecking}.
However, general-purpose model checkers such as
SPIN~\citep{holzmann1997spin} and NuSMV~\citep{cimatti2002nusmv}
require manual translation of the system under analysis into a
dedicated modeling language---a prohibitive overhead for agent
developers iterating on workflow designs. Similarly, the business
process management (BPM) community has a long tradition of workflow
soundness verification~\citep{vanderaalst2011bpm}, but these methods
target visual modeling tools and standardized notations (e.g., BPMN,
Petri nets), not the programmatic API objects through which agent
frameworks define their graphs. To date, no system automatically
extracts and statically verifies agent workflow graphs directly from
framework source code.

\paragraph{Our approach.}
This paper argues for \emph{pre-deployment static verification} of
agent workflow graphs. Our key observation is that modern agent
frameworks expose workflow structure through their APIs, making the
graph amenable to automated extraction and analysis. When agent behavior
is graph-structured, safety properties reduce to reachability,
isolation, and temporal constraints over the topology---properties that
can be checked exhaustively without running the
agent~\citep{clarke1999modelchecking}.

We present \textsc{Agentproof}, a practical system for static analysis
of agent workflows that focuses on properties (a)~expressible over the
workflow graph and (b)~independent of LLM text generation semantics.
Unlike general-purpose model checkers, Agentproof extracts the analysis
model automatically from framework source code, requiring no manual
modeling effort.

\paragraph{Contributions.}
\begin{itemize}
  \item A cross-framework extraction and integration pipeline: we show
  that four major agent frameworks (LangGraph, CrewAI, AutoGen, Google
  ADK) expose enough structure through their APIs to enable fully
  automatic extraction into a \emph{unified abstract workflow model}
  with typed nodes and edges---eliminating the manual modeling step
  required by general-purpose model
  checkers~(Section~\ref{sec:graph_model}).
  \item Extractors for LangGraph, CrewAI, AutoGen, and Google ADK that
  bridge heterogeneous framework representations---each with different
  entry/exit conventions, edge semantics, and hierarchy
  models---into a single \texttt{AgentGraph}.
  \item A pre-deployment verification pipeline comprising five
  structural checks with \emph{witness trace generation} and a temporal
  policy DSL covering the safety fragment of LTL, compiled to
  deterministic finite automata~(Section~\ref{sec:verification_methods}).
  \item An empirical evaluation on 18 curated workflows demonstrating
  that structural defects (dead-end nodes, unreachable exits) and
  policy violations (missing human gates) are common, with sub-second
  verification for graphs up to 5{,}000
  nodes~(Section~\ref{sec:evaluation}).
\end{itemize}

\paragraph{Scope.}
Our approach validates workflow design decisions early: ``can this
graph reach a destructive tool?'' or ``must human review occur between
two sensitive actions?'' It does not attempt to prove semantic
properties of LLM outputs (e.g., truthfulness), which remain
inherently non-deterministic. We discuss this boundary precisely in
our threat model (Section~\ref{sec:threat_model}).
